\section{Derivace}
\subsection{Definice a základní vztahy}
\begin{definice}
Nechť $a\in\R$ a $f$ je funkce. Existuje-li limita
\[
\lim_{h\to 0} \frac{f(a+h)-f(a)}{h},
\]
pak její hodnotu nazýváme \textbf{derivací} funkce $f$ v bodě $a$. Obdobně definujeme \textbf{derivaci zprava} a \textbf{derivaci zleva} funkce $f$ v bodě $a$ předpisy
\[
\lim_{h\to 0^+} \frac{f(a+h)-f(a)}{h},\qquad \lim_{h\to 0^-} \frac{f(a+h)-f(a)}{h}.
\]
Derivaci zleva v bodě a derivaci zprava v bodě nazýváme \textbf{jednostrannými derivacemi}.
\end{definice}

\begin{oznaceni}
	Pokud derivace existuje, je určena jednoznačně (důsledek věty \ref{3.1}). Používáme značení
	$f^\prime(a)$, $f_+^\prime(a)$, $f_-^\prime(a)$.
\end{oznaceni}

\begin{poznamka}\,
\begin{enumerate}[$i$.]
	\item Derivace
	$$f^\prime(a)
 \begin{cases}
		\textrm{existuje}, \\
		\textrm{neexistuje} \,\,\,\begin{cases}
			\textrm{je vlastní, tj. }\in \mathbb R, \\
			\textrm{je nevlastní} \,\,\, \begin{cases}
				\infty, \\
				-\infty.
\end{cases}
\end{cases}
 \end{cases}
	$$
	\item $f^\prime (a) \implies \exists \delta \in \mathbb R , \delta > 0: f$ je definována na okolí $B(a,\delta)$
	\item $f^\prime(a) = A\in \mathbb R ^* \iff f_+^\prime(a) = f_-^\prime (a) = A\in \mathbb R ^*$
	\item $\lim_{x\to a} \frac{f(x)-f(a)}{x-a} = \lim_{h\to 0} \frac{f(a+h)-f(a)}{h}$
	\item Derivace je lokální pojem.
\end{enumerate}
\end{poznamka}

\begin{poznamka}[Derivace a tečna]
	Nechť $a\in \mathbb R $, $f$ je funkce a $f^\prime(a)\in \mathbb R $. Tečnou ke grafu $f$ v bodě
	$\left [ a,f(a) \right ] $ rozumíme funkci
	$$
	   t: x\mapsto f(a)+ f^\prime(a)(x-a).
	$$
	Platí
	\begin{align*}
	\lim_{x\to a} \frac{f(x)-t(x)}{x-a}&= \lim_{x\to a} \frac{f(x)-f(a)-f^\prime(a)(x-a)}{x-a}\\
	&= \left( \frac{f(x)-f(a)}{x-a}-f(a) \right )= f^\prime(a)-f^\prime(a) = 0.
	\end{align*}
\end{poznamka}

\begin{poznamka}
	Definičním oborem derivace funkce $f$ rozumíme množinu
	$\left \{ x \in \mathbb R , f^\prime (x)\in \mathbb R  \right \} $.
\end{poznamka}

\begin{priklad}\,
\begin{enumerate}[$i.$]
	\item $f(x)=c, c\in B(a, \delta)$, kde $\delta>0, a \in \mathbb R $. Platí
	$$
	    f^\prime(a) = \lim_{x\to a} \frac{f(x)-f(a)}{x-a} = \lim_{x\to a} \frac{c-c}{x-a} = 0.
	$$
	\item $f(x)=x^m$, kde $m\in\mathbb{N}$, $x\in\mathbb{R}$:
	\par  Pro $a\in\mathbb{R}$:
    \[
    f'(a)=\lim_{x\to a}\frac{x^m-a^m}{x-a}=\lim_{x\to a}\frac{(x-a)\left(x^{m-1}+x^{m-2}a+\cdots+a^{m-1}\right)}{x-a}=na^{n-1}.
    \]
    \item $f(x)=\operatorname{sgn}x$, $x\in\mathbb{R}$
    \begin{align*}
    f'_+(0)&=\lim_{x\to0^+}\frac{\operatorname{sgn}x-\operatorname{sgn}0}{x-0}
    =\lim_{x\to0^+}\frac{1}{x}=+\infty\\
    f'_-(0)&=\lim_{x\to0^-}\frac{\operatorname{sgn}x-\operatorname{sgn}0}{x-0}
    =\lim_{x\to0^-}\frac{-1}{x}=+\infty,
    \end{align*}
    odkud máme, že $\operatorname{sgn}^\prime(0)=\infty.$
    \item $f(x)=|x|$
    \[
    f'_+(0)=1, \qquad f'_-(0)=-1,
    \]
    a proto derivace $f'(0)$ neexistuje.
\end{enumerate}
\end{priklad}

\begin{veta}[Derivace a spojitost]\label{veta5.1}
	Nechť funkce $f$ má v bodě $a \in \mathbb R $ vlastní derivaci. Potom je $f$ v bodě $a$ spojitá.
\end{veta}

\begin{proof}
	Počítejme
    $$
	\lim_{x\to a} f(x) = \lim_{x\to a} \left( \frac{f(x)-f(a)}{x-a}(x-a)+f(a) \right ) =\underbrace{f^\prime(a)}_{\in \mathbb R }\cdot 0+f(a)=f(a).
	$$
\end{proof}

\begin{poznamka}\,
\begin{enumerate}[(i)]
	\item Věta \ref{veta5.1} platí i v jednostranné verzi.
	\item Předpoklad, že $f^\prime(a)$ je vlastní, nelze vynechat (vizte funkci signum v nule).
\end{enumerate}
\end{poznamka}

\begin{veta}[Aritmetika derivací]\label{veta5.2}
Nechť $a \in \mathbb{R}$ a $f, g$ jsou funkce, které mají v bodě $a$ derivaci.
\begin{itemize}
\item[(a)] Pak platí $$(f+g)'(a)=f'(a)+g'(a),$$ je-li výraz na pravé straně definován.
\item[(b)] Nechť alespoň jedna z funkcí $f$ a $g$ je spojitá v bodě $a$.
Potom platí $$(fg)'(a)=f'(a)g(a)+f(a)g'(a),$$
je-li výraz na pravé straně definován.
\item[(c)] Nechť funkce $g$ je spojitá v bodě $a$. Potom platí
$$
\left(\frac{f}{g}\right)'(a)=\frac{f'(a)g(a)-f(a)g'(a)}{g^2(a)},$$
je-li výraz na pravé straně definován.
\end{itemize}
\end{veta}

\begin{proof}\,
\begin{enumerate}[(a)]
	\item Platí
	\begin{align*}
	(f+g)'(a) & = \lim_{h\to 0} \frac{(f+g)(a+h)-(f+g)(a)}{h}\\
	&= \lim_{h\to 0} \frac{f(a+h)+g(a+h)-f(a)-g(a)}{h}\\
	&= \lim_{h\to 0} \left( \frac{f(a+h)-f(a)}{h} + \frac{g(a+h)-g(a)}{h} \right)\\
 &= f'(a) + g'(a).
\end{align*}
\item Předpokládejme, že $g$ je spojitá v $a$. Pak
\begin{align*}
(fg)'(a) &= \lim_{h\to 0} \frac{(fg)(a+h)-(fg)(a)}{h}\\
&= \lim_{h\to 0} \frac{f(a+h)g(a+h)-f(a)g(a)}{h}\\
&= \lim_{h\to 0} \frac{(f(a+h)-f(a))g(a+h)+f(a)(g(a+h)-g(a))}{h}\\
&= f'(a)\,g(a) + f(a)\,g'(a).
\end{align*}
\item Platí
\begin{align*}
\left(\frac{f}{g}\right)'(a) &= \lim_{h\to 0} \frac{\left( \frac{f}{g} \right ) (a+h)-\left( \frac{f}{g} \right ) (a)}{h}\\
	&= \lim_{h\to 0} \frac{1}{h}\left( \frac{f(a+h)}{g(a+h)} - \frac{f(a)}{g(a)} \right)\\
	&= \lim_{h\to 0} \frac{1}{h}\,\frac{f(a+h)g(a)-f(a)g(a+h)}{g(a+h)g(a)}\\
	&= \lim_{h\to 0} \frac{1}{g(a+h)g(a)}\left( \frac{f(a+h)-f(a)}{h}g(a) - f(a)\frac{g(a+h)-g(a)}{h} \right)\\
	&= \frac{f'(a)g(a)-f(a)g'(a)}{g(a)^2}.\qedhere
\end{align*}
\end{enumerate}
\end{proof}

\begin{veta}[Derivace složené funkce] \label{veta5.3}
	Nechť funkce $g$ je spojitá v bodě $a \in \mathbb R$, má v tomto bodě derivaci
	a funkce $f$ má derivaci v bodě $g(a)$. Pak platí
	$$
	    (f\circ g)^\prime(a) = f^\prime(g(a))\cdot g^\prime(a),
	$$
	je-li výraz na pravé straně definován.
\end{veta}

\begin{proof}
	Označme $b=g(a)$,
	$$
	\varphi(y)=\frac{f(y)-f(b)}{y-b}, \quad \mathcal D(f)\smallsetminus \left \{ b \right \} .
	$$
	Platí $\lim_{y\to b} \varphi(y)=f^\prime(b)$. Rozlišíme dva případy
	\begin{enumerate}[(i)]
		\item $g^\prime (a) \ne 0$: nalezneme $\eta > 0$ takové, že platí
		$$
		    \forall x in P(a,\eta):\frac{g(x)-g(a)}{x-a} \ne 0.
		$$
		odtud plyne, že
		$$
		    \forall x \in P(a, \eta): g(x)\ne g(a).
		$$
		Počítejme
		\begin{align*}
			(f\circ g)^\prime (a) &= \lim_{x\to a} \frac{f(g(x))-f(g(a))}{x-a}=\lim_{x\to a} \frac{f(g(x))-f(g(a))}{g(x)-g(a)}\cdot \frac{g(x)-g(a)}{x-a}\\
			&=\lim_{x\to a} \varphi(g(x))\cdot \frac{g(x)-g(a)}{x-a}.
		\end{align*}
		\item $g^\prime(a) = 0$: Výraz na pravé straně je definován, a proto
		$f^\prime(b)\in \mathbb R $. Nalezneme $c>0, \delta>0$ taková, že
		$$
		    \forall y \in P(b,\delta): |\varphi(y)|<c \quad (\textrm{věta }\ref{veta3.2}).
		$$
		Dále
		\begin{align*}
			& \forall y \in P(b,\delta): \left | \frac{f(y)-f(b)}{y-b} \right |<c,\\
			& \forall y \in B(b,\delta): |f(y)-f(b)| \leq c\cdot |y-b|.
\end{align*}
    Zvolme $\varepsilon>0$. Nalezneme $\rho_1>0$ takové, že
    $$
        \forall x \in P(a,\rho_1): \left | \frac{g(x)-g(a)}{x-a}-0 \right|<\varepsilon.
    $$
    Dále nalezneme $\rho_2 \in (0,\rho_1)$ takové, že
    $$
        \forall x \in B(a,\rho_2): g(x)\in B(b,\delta).
    $$
    Pro každé $x \in P(a,\rho_2)$ platí
    $$
	\left | \frac{f(g(x))-f(g(a))}{x-a} \right | \leq \frac{c|g(x)-g(a)|}{|x-a|} = c\cdot \left | \frac{g(x)-g(a)}{x-a} \right |<c\varepsilon,
    $$
    odkud plyne $(f\circ g)^\prime(a)=0=f^\prime(b)\cdot 0 = f^\prime(b)\cdot g^\prime(a)$.\qedhere
\end{enumerate}
\end{proof}

\begin{poznamka}
	V čem je problém:
	\begin{align*}
    (f\circ g)^\prime (a) &= \lim_{x\to a} \frac{f(g(x))-f(g(a))}{x-a} = \lim_{x\to a} \frac{f(g(x))-f(g(a))}{g(x)-g(a)}\cdot \frac{g(x)-g(a)}{x-a}\\
    &= f^\prime(b)\cdot g^\prime(a).
    \end{align*}
    Funkce
    \begin{align*}
	&y\mapsto \frac{f(y)-f(b)}{y-b} && \lim_{y\to b} \frac{f(y)-f(b)}{y-b}=f^\prime(b)-f^\prime(g(a))\\
	&g(a)=b && \lim_{x\to a} g(x) = g(a)=b
    \end{align*}
\end{poznamka}

\begin{veta}[Derivace inverzní funkce] \label{veta5.4}
Nechť $a$ je vnitřním bodem intervalu $I$ a $f$ je spojitá a ryze monotónní funkce na $I$. Označme $b=f(a)$.
\begin{enumerate}[(a)]
  \item Nechť $f'(a)$ existuje a je nenulová. Potom platí
  $$(f^{-1})'(b)=\dfrac{1}{f'(a)}.$$
  (Je-li $f'(a)=\pm\infty$, pak $(f^{-1})'(b)=0$.)
  \item Nechť $f'(a)=0$ a $f$ je rostoucí na $I$. Potom $(f^{-1})'(b)=\infty$.
  \item Nechť $f'(a)=0$ a $f$ je klesající na $I$. Potom $(f^{-1})'(b)=-\infty$.
\end{enumerate}
\end{veta}

\begin{proof}
	Funkce $f:I\to \mathbb R $. Označme $J=f(I)$. Podle věty \ref{veta3.12} je $J$ interval.
	Inverzní funkce $f^{-1}$ je tedy $J\to I$. Podle věty \ref{veta3.15} je $f^{-1}$ spojitá
	a ryze monotónní. Označme
	$$
	    \varphi(x) = \frac{f(x)-f(a)}{x-a}, \quad x \in \mathcal D(f)\smallsetminus \left \{ a \right \}.
	$$
	Platí $\lim_{x\to a} \varphi(x)=f^\prime(a)$ v případech (a) až (c). Dále platí
	$\lim_{y\to b} f^{-1}(y)=f^{-1}(b)=a$, neboť $b$ je vnitřním bodem $J$ a $f^{-1}$ je spojitá
	na intervalu $J$.
	\begin{enumerate}[(a)]
		\item $\lim_{y\to b} (\varphi \circ f^{-1})(y)$, a je splněna podmínka (P), protože
		$f^{-1}$ je ryze monotónní. Máme tedy
		$$
  \lim_{y\to b} \frac{f(f^{-1}(y))-f(a)}{f^{-1}(y)-a} = \lim_{y\to b} \frac{y-b}{f^{-1}(y)-f^{-1}(b)},
		$$
		tedy
		$$
		    (f^{-1})^\prime(y) = \lim_{y\to b} \frac{f^{-1}-f^{-1}(b)}{y-b} = \lim_{y\to b} \frac{1}{\frac{y-b}{f^\prime(y)-f^\prime(b)}}=\frac{1}{f^\prime(a)}.
		$$
		\item Podobně jako v (a) máme, že
		$$\lim_{y\to b} \frac{y-b}{f^{-1}(y)-f^{-1}(b)}=f^\prime (a) = 0.$$
		Potom platí podle věty \ref{veta3.4}
		$$
		    (f^{-1})^\prime = \lim_{y\to b} \frac{1}{\frac{y-b}{f^{-1}(y)-f^{-1}(b)}}=\infty.
		$$
		\item Analogicky jako (b).\qedhere
	\end{enumerate}
\end{proof}

\paragraph{Derivace elementárních funkcí}
\begin{itemize}
	\item $f(x)=c, x \in \mathbb R$
	\par $f^\prime(x) = 0, x \in \mathbb R $
	\item $\left( x^n \right )^\prime = nx^{n-1}, x \in \mathbb R , n \in \mathbb N  $
	\item $\left( \exp(x) \right )^\prime  = \exp(x), x \in \mathbb R $
	$$
 \left( \exp(x) \right )^\prime = \lim_{h\to 0} \frac{\exp(x+h)-\exp(x)}{h} = \lim_{h\to 0} \exp (x)\cdot \frac{\exp(h)-1}{h} = \exp(x)
	$$
	\item $(\log x)^\prime = \frac{1}{x}, x \in \left( 0,\infty \right ) $
	$$
	    (\log x)^\prime =\frac{1}{\exp(\log x)} = \frac{1}{x}.
	$$
	\begin{align*}
		f(z) &= \exp(z) \\
		f^{-1}(x) & = \log x
	\end{align*}
	\item $(x^a)^\prime = ax^{a-1}, x \in \left( 0,\infty \right ), a \in \mathbb R  $
	$$
 \left( x^a \right )^\prime = \left( e^{a \log x} \right )^\prime = e^{a\log x}\cdot \frac{a}{x} = ax^{a-1}.
	$$
	\item $\left( a^x \right )^\prime = a^x \log a, x \in \mathbb R , a > 0 $
	$$
 \left( a^x \right )^\prime = \left( e^{x\log a} \right )^\prime = e^{x\log a}\cdot \log a = a^x\log a.
	$$
\end{itemize}

\begin{poznamka}
	Funkce $x^x$ se nehodí derivovat podle žádného z předchozích bodů, ale
$$\left( x^x \right )^\prime = \left( e^{x\log x} \right )^\prime = e^{x\log x}\cdot (x\log x + x\cdot \frac{1}{x}) = x^x (\log x +1), x > 0.  $$
\end{poznamka}

\begin{itemize}
	\item $\left( \sqrt[n]{x} \right )^\prime = \frac{1}{n} \cdot \frac{\sqrt[n]{x}}{x} $ pro $x \in (0,\infty)$ pro $n \in \mathbb N $ sudé, nebo $x \in \mathbb R \smallsetminus \{0\}$ pro $n \in \mathbb N $ liché.
	Dále $(\sqrt[n]{x})^\prime = \infty$ pro $x=0$ a $n \in \mathbb N $ liché,
	$(\sqrt[n]{x})_+^\prime = \infty$ pro $x=0$ a $n\in \mathbb N $ sudé.
	\item $(\sin x )^\prime = \cos x, x \in \mathbb R $
	\begin{align*}
	(\sin x )^\prime &=\lim_{h\to 0} \frac{\sin(x+k)-\sin x}{h} = \lim_{h\to 0} \frac{\sin x \cos h + \cos x \sin h-\sin x}{h} \\
	&= \lim_{h\to 0} \left( \sin x \cdot \frac{\cos h - 1}{h} + \cos x \cdot \frac{\sin h }{h} \right ),
	\end{align*}
	ale
	$$
 \lim_{h\to 0} \frac{\cos h -1}{h} = \lim_{h\to 0} \frac{\cos h-1}{h^2}\cdot h = \left( -\frac{1}{2} \right )\cdot 0,
	$$
	tedy $(\sin x )^\prime = \cos x.$
	\item $(\cos x)^\prime = -\sin x \, x \in \mathbb R $
	\item $(\tg x)^\prime = \frac{1}{\cos^2 x}, x \in \mathcal D(\tg)$
	\item $(\cotg x )^\prime = -\frac{1}{\sin^2 x}, x \in \mathcal D(\cotg)$
	$$
	    (\cotg x) ^\prime = \left( \frac{\cos x}{\sin x} \right )^\prime = \frac{-\sin x \cdot \sin x - \cos x \cos x}{\sin^2 x} = - \frac{1}{\sin^2 x}.
	$$
	\item $(\arcsin x)^\prime = \frac{1}{\sqrt{1-x^2}}, x \in (-1,1)$
	\begin{align*}
		f(z) &= \sin z, z \in \left [ -\frac{\pi}{2}, \frac{\pi}{2} \right ]  \\
		f^{-1}(x) &= \arcsin x, x \in \left [ -1,1 \right ] \\
		(\arcsin x)^\prime &= \frac{1}{\cos(\arcsin x)}, \quad x \in (-1,1)\\
		\cos^2 y &= 1-\sin^2 y, \quad \ in \left( -\frac{\pi}{2}, \frac{\pi}{2} \right ) \implies \cos y = \sqrt{1-\sin ^2y} \\
		\cos(\arcsin x) &= \sqrt{1-\sin^2(\arcsin x)}=\sqrt{1-x^2}.
	\end{align*}
	\item $(\arccos x)^\prime = -\frac{1}{\sqrt{1-x^2}}, x \in (-1,1)$
	\item $(\arctg x)^\prime = \frac{1}{1+x^2},x \in \mathbb R $
	\item $(\arccotg x)^\prime = -\frac{1}{1+x^2}, x \in \mathbb R $
\end{itemize}

\begin{veta}[Nutná podmínka lokálního extrému] \label{veta5.5}
Jestliže $a$ je bodem lokálního extrému funkce $f$, potom buď $f'(a)$ neexistuje, nebo $f'(a)=0$.
\end{veta}

\begin{proof}
	Pro spor předpokládejme, že $f^\prime(a)$ existuje a $f^\prime(a)\ne 0$. Předpokládejme nejprve, že
	$f^\prime(a)>0$. Potom
	$$
	    f^\prime(a) = \lim_{x\to a} \frac{f(x)-f(a)}{x-a} > 0.
	$$
	Nalezneme $\delta > 0$ takové, že
	$$
	    \forall x \in P(a, \delta) : \frac{f(x)-f(a)}{x-a} > 0.
	$$
	Potom
	\begin{align*}
	    \forall x \in (a,a+\delta): f(x)>f(a), && \forall x \in (a-\delta, a): f(x) < f(a).
	\end{align*}
	Proto $a$ není bodem lokálního extrému, což je spor. Případ $f^\prime(a)<0$ lze přivést
	ke sporu analogicky.
\end{proof}

\begin{priklad}[Častá úloha]
	Nechť $f:\left [ a,b \right ] \to \mathbb R $ je spojitá na $\left [ a,b \right ] $ a
	pro každé $x \in (a,b)$ existuje $f^\prime(x)$. Nalezněte extrémy $f$ na $\left [ a,b \right ] $.
\end{priklad}

\begin{reseni}
	Protože je $f$ spojitá na $\left [ a,b \right ] $, má $f$ maximum i minimum.
	Dále ať $$E=\left \{ x\in (a,b), f^\prime(x)=0 \right \} \cup \left \{ a,b \right \}.  $$
	Extrémy mohou být jedině v $E$. Dále vypočteme funkční hodnoty prvků v $E$ a podle toho
	určíme maximum a minimum.
\end{reseni}

\subsection{Věty o střední hodnotě}
\begin{veta}[Rolle] \label{veta5.6}
Nechť $a,b \in \R, a<b$ a $f\colon[a,b]\to\R$ je spojitá funkce. Je-li $f(a)=f(b)$ a $f$ má derivaci v každém bodě intervalu $(a,b)$, pak existuje $c\in(a,b)$ takové, že $f'(c)=0$.
\end{veta}

\begin{proof}
	Podle věty \ref{veta3.14} $f$ nabývá v intervalu $\left [ a,b \right ] $ svého maxima i minima.
	Označme $m=\min f (\left [ a,b \right ])  $ a $M = \max f (\left [ a,b \right ])$.  Potom platí
	$$
	    m \leq f(a) = f(b) \leq M.
	$$
	\par Pokud $m=M$, pak je $f$ konstantní na $\left [ a,b \right ] $ a stačí volit $c \in (a,b)$ libovolně.
	\par Pokud $m < M$, pak platí $m< f(a)$ nebo $f(a)< M$. Pokud $f(a)<M$, pak existuje $c\in \left [ a,b \right ] $
	takové, že $f(c)=M.$ Podle předpokladů věty a podle věty \ref{veta5.5} máme $f^\prime(c)=0$.
	Pokud $m<f(a)$, postupujeme obdobně.
\end{proof}

\begin{veta}[Lagrange] \label{veta5.7}
Nechť $a,b \in \R, a<b$ a $f\colon[a,b]\to\R$ je spojitá funkce, která má v každém bodě intervalu $(a,b)$ derivaci. Pak existuje $c\in(a,b)$ takové, že
\[
 f'(c)=\frac{f(b)-f(a)}{b-a}.
\]
\end{veta}

\begin{proof}
Definujme pomocnou funkci
$$
    g(x)=f(x) - \frac{f(b)-f(a)}{b-a}(c-a), \quad x \in \left [ a,b \right ] .
$$
Platí, že
\begin{itemize}
	\item $g$ je spojitá na $\left [ a,b \right ] $,
	\item $g(a)= f(a)$, $g(b)=f(b)-\frac{f(b)-f(a)}{b-a}(b-a)=f(a)$,
\item $g^\prime(x)=\underbrace{f^\prime(x)}_{\in \mathbb R ^*}-\underbrace{\frac{f(b)-f(a)}{b-a}}_{\in \mathbb R }$, $x \in (a,b)$.
\end{itemize}
Podle Rolleovy věty \ref{veta5.6} nalezneme $c\in(a,b)$ takové, že $g^\prime(c)=0$. Potom platí
$$
    \underbrace{f^\prime(c)}_{\in \mathbb R }-\frac{f(b)-f(a)}{b-a}=0.
$$
Odsud dostáváme
$$
    f^\prime(c)=\frac{f(b)-f(a)}{b-a},
$$
což jsme chtěli dokázat.
\end{proof}

\begin{veta}[Cauchy] \label{veta5.8}
Nechť $a,b \in \R, a<b$, $f$ a $g$ jsou funkce spojité na $[a,b]$, $f$ má v každém bodě $x\in(a,b)$ derivaci a $g$ má v každém bodě $x\in(a,b)$ vlastní nenulovou derivaci. Potom $g(a)\ne g(b)$ a existuje $c\in(a,b)$ takové, že
\[
 \frac{f'(c)}{g'(c)}=\frac{f(b)-f(a)}{g(b)-g(a)}.
\]
\end{veta}

\begin{proof}
	Podle Lagrangeovy věty \ref{veta5.7} nalezneme $d \in (a,b)$ takový, že
	$$
	    g^\prime(d)=\frac{g(b)-g(a)}{b-a}.
	$$
	Platí $g^\prime(d)\ne 0$, a tedy $g(a)\ne g(b)$. Položme
	$$
	\varphi(x) = (f(x)-f(a))(g(b)-g(a))-(f(b)-f(a))(g(x)-g(a)), \quad x \in \left [ a,b \right ].
	$$
	Platí, že
	\begin{itemize}
		\item $\varphi$ je spojitá na $\left [ a,b \right ] $,
		\item $\varphi(a)=0$, $\varphi(b)=0$,
		\item $\varphi^\prime(x)=\underbrace{f^\prime(x)}_{\in \mathbb R ^*}\underbrace{(g(b)-g(a))}_{\ne 0}-(f(b)-f(a))\underbrace{g^\prime(x)}_{\in \mathbb R}$, $x \in (a,b)$.
	\end{itemize}
	Podle Rolleovy věty \ref{veta5.6} nalezneme $c \in (a,b)$ takové, že $\varphi^\prime(c)=0$,
	neboli
	$$
	    \underbrace{f^\prime(c)}_{\in \mathbb R }\underbrace{(g(b)-g(a))}_{\ne 0}-(f(b)-f(a))\underbrace{g^\prime(c)}_{\ne 0}=0.
	$$
	Odtud plyne
	$$
 \frac{f^\prime(c)}{g^\prime(c)}=\frac{f(b)-f(a)}{g(b)-g(a)},
	$$
	což jsme chtěli dokázat.
\end{proof}

\begin{veta}[l'H\^ospitalovo pravidlo] \label{veta5.9}
Nechť $a\in\R\cup\{-\infty\}$, $f,g$ jsou funkce, existuje
$$\displaystyle \lim_{x\to a^+}\frac{f'(x)}{g'(x)}$$ a platí buď
\begin{enumerate}[(a)]
  \item $\lim_{x\to a^+} f(x)=\lim_{x\to a^+} g(x)=0$, nebo
  \item $\lim_{x\to a^+} \abs{g(x)}=\infty$.
\end{enumerate}
Potom platí $$ \lim_{x\to a^+}\frac{f(x)}{g(x)}=\lim_{x\to a^+}\frac{f'(x)}{g'(x)}.$$
\end{veta}

\begin{proof}
	Provedeme důkaz (a), důkaz (b) umět nemusíme.
	\par Označme $L=\lim_{x\to a_+} \frac{f^\prime(x)}{g^\prime(x)}.$
	\paragraph{Pomocné tvrzení.}
	\begin{enumerate}
		\item[(i)] Pro každé $\alpha \in \mathbb R ,\alpha > L$ existuje $\delta > 0$ takové, že
		$$
		    \forall x \in P_+(a,\delta):\frac{f(x)}{g(x)}<\alpha.
		$$
		\item[(ii)] Pro každé $\beta \in \mathbb R , \beta < L$ existuje $\delta^\prime>0$ takové,
		že
		$$
		    \forall x \in P_+(a,\delta^\prime):\frac{f(x)}{g(x)}>\beta.
		$$
	\end{enumerate}
		\paragraph{Důkaz pomocného tvrzení.}
		\begin{enumerate}
			\item[(i)] Nalezneme $r\in (L,\alpha)$. Nalezneme $\delta > 0$ takové, že platí
			\begin{itemize}
				\item $f$, $f^\prime$, $g$ a $g^\prime$ jsou definovány na $P_+(a,\delta)$,
				\item $f^\prime$, $g^\prime$ jsou vlastní na $P_+(a,\delta)$,
				\item $g^\prime$ je nenulové na $P_+(a,\delta)$,
				\item $\forall x \in P_+(a,\delta):\frac{f^\prime(x)}{g^\prime(x)}<r$,
				\item $g$ je nenulová na $P_+(a,\delta)$; pokud by existoval nulový bod $f$ v
				$P_+(a,\delta)$, pak je nejvýše jeden a $\delta$ vhodně zmenšíme.
			\end{itemize}
		Zvolme $x\in P_+(a,\delta)$. Předpokládejme, že $y \in P_+(a,\delta)$ a $y<x$. Podle Cauchyovy
		věty \ref{veta5.8} nalezneme $c \in (y,x)$ takové, že
		$$
  \frac{f^\prime(c)}{g^\prime(c)}=\frac{f(x)-f(y)}{g(x)-g(y)}.
		$$
		Potom $\frac{f^\prime(c)}{g^\prime(c)}<r$, a tedy
		$$
  \frac{f(x)-f(y)}{g(x)-g(y)}<r.
		$$
		Potom máme, že
		$$
  \lim_{y\to a^+} \frac{f(x)-f(y)}{\underbrace{g(x)-g(y)}_{<r}}=\underbrace{\frac{f(x)}{\underbrace{g(x)}_{\ne 0}}}_{\leq r}<\alpha.
		$$
		\item[(ii)] Dokážeme analogicky.
	\end{enumerate}
	\paragraph{Vlastní důkaz.} Pokud $L= \infty$, resp. $L=-\infty$, pak věta plyne z (ii), resp. z (i).
	Pokud $L\in \mathbb R $ a $\varepsilon > 0$, položíme $\alpha = L+\varepsilon$, $\beta = L-\varepsilon$
	a nalezneme příslušná $\delta$ a $\delta^\prime$. Položme $\delta_1=\min \left \{ \delta, \delta^\prime \right \} $.
	Pak pro každé $x \in P_+(a,\delta_1)$ platí
	$$
	    L-\varepsilon =\beta < \frac{f(x)}{g(x)}<\alpha = L+\varepsilon, \quad \textrm{tj. } \frac{f(x)}{g(x)}\in B(L,\varepsilon),
	$$
	což jsme chtěli dokázat.
\end{proof}

\begin{poznamka}\,
\begin{enumerate}[$i.$]
	\item $\lim_{x\to 0} \frac{2x+1}{3x+1} \ne \lim_{x\to 0} \frac{2}{3}$
	\item $\lim_{x\to \infty} \frac{x}{x+\sin x}=1$, ale l'H\^ospitalovým pravidlem výsledek nedostaneme
	\item $\lim_{x\to 0} \frac{e^{-\frac{1}{x^2}}}{x}= \lim_{x\to 0} \frac{e^{-\frac{1}{x^2}}\cdot \frac{2}{x^3}}{1}=\lim_{x\to 0} \frac{2e^{-\frac{1}{x^2}}}{x^3}$
\end{enumerate}
\end{poznamka}

\begin{priklad}
	Spočtěte $\lim_{x\to 0} \frac{\sin x -x}{x^3}$.
\end{priklad}

\begin{reseni}
	L'H\^ospitalovým pravidlem dostáváme
	$$
	    \lim_{x\to 0} \frac{\sin x -x}{x^3} = \lim_{x\to 0} \frac{\cos x -1}{3x^2} = \frac{1}{3}\cdot \frac{-1}{2} = -\frac{1}{6}.
	$$
\end{reseni}

\begin{veta}[Vztah derivace a monotonie]\label{veta5.10}
Nechť $f$ je spojitá na intervalu $J$ a v každém vnitřním bodě $J$ (množinu všech vnitřních bodů intervalu $J$ označme jako $\operatorname{int} J$) existuje derivace $f$.
\begin{enumerate}[(i)]
  \item Je-li $f'(x)>0$ pro všechna $x\in\operatorname{int} J$, pak $f$ je rostoucí na $J$.
  \item Je-li $f'(x)<0$ pro všechna $x\in\operatorname{int} J$, pak $f$ je klesající na $J$.
  \item Je-li $f'(x)\ge 0$ pro všechna $x\in\operatorname{int} J$, pak $f$ je neklesající na $J$.
  \item Je-li $f'(x)\le 0$ pro všechna $x\in\operatorname{int} J$, pak $f$ je nerostoucí na $J$.
\end{enumerate}
\end{veta}

\begin{proof}
Dokážeme (i), ostatní případy se dokážou analogicky.
	Předpokládejme, že $x,y\in J$, $x<y$. Podle Lagrangeovy věty \ref{veta5.7} nalezneme
	$c \in (x,y)$ takové, že
	$$
	    f^\prime(c) = \frac{f(y)-f(x)}{y-x}>0.
	$$
	Odtud plyne $f(x)<f(y)$.
\end{proof}

\begin{poznamka}[Varování!]
	Věta \ref{veta5.10} je formulovaná pro \emph{interval}.
\end{poznamka}

\begin{veta}\label{veta5.11}
Nechť $f$ je spojitá zprava v bodě $a\in\R$ a existuje $\displaystyle \lim_{x\to a^+} f'(x)$. Potom existuje $f'_+(a)$ a platí $f'_+(a)=\displaystyle \lim_{x\to a^+} f'(x)$.
\end{veta}

\begin{proof}
    Je vlastně
	$$
	f_+^\prime(a)=\lim_{x\to a_+} \frac{f(x)-f(a)}{x-a}.
	$$
	Platí, že
	$$
 \lim_{x\to a_+} (f(x)-f(a))=0, \qquad \lim_{x\to a_+} (x-a)=0.
	$$
	Odtud z l'H\^ospitalova pravidla dostáváme
	$$
	    f_+^\prime(a) = \lim_{x\to a_+} \frac{f^\prime(x)}{1}=\lim_{x\to a_+} f^\prime(x).
	$$
\end{proof}

\begin{definice}
Nechť $n\in\N$, $a\in\R$ a $f$ má vlastní $n$-tou derivaci na okolí bodu $a$. Potom \textbf{$(n+1)$-ní derivací funkce} $f$ v bodě $a$ rozumíme
\[
 f^{(n+1)}(a)=\lim_{h\to 0} \frac{f^{(n)}(a+h)-f^{(n)}(a)}{h}.
\]
\end{definice}

\begin{poznamka}
	Derivace značíme
	$$
	    f^\prime \quad f^{\prime\prime}\quad f^{\prime\prime\prime}\quad f^{(4)}\quad f^{(5)}\quad\dots\quad f^{(n)}.
	$$
\end{poznamka}
